\documentclass[10pt,a4paper]{article}
\usepackage[T1]{fontenc}
\usepackage[utf8]{inputenc}
\usepackage{lmodern,microtype}
\usepackage[margin=18mm]{geometry}
\usepackage{amsmath,amssymb,booktabs,tabularx,array}
\usepackage{xcolor,hyperref,enumitem}
\definecolor{navy}{HTML}{17365D}
\hypersetup{colorlinks=true,linkcolor=navy,urlcolor=navy}
\setlength{\parindent}{0pt}
\setlength{\parskip}{0.55\baselineskip}
\setlist{nosep,leftmargin=*}
\newcolumntype{Y}{>{\raggedright\arraybackslash}X}
\newcommand{\code}[1]{\texttt{\detokenize{#1}}}
\title{\textbf{Curriculum Learning for Multi-user Forecasting}\\Experiment Guideline}
\author{Thesis project}
\date{25 August 2026}
\begin{document}
\maketitle

\section{Research setting}
The project tests whether ordering users from easier to harder improves
forecasting convergence, final accuracy, or tail-user performance.  A run holds
the forecasting architecture, optimizer-step budget, batch size,
normalization, loss, chronological splits, and evaluation windows fixed; only
the user-sampling curriculum changes.

For user $u$, a sampled query date $t$ supplies
$X_{u,t}=X_u(t-L,t]$ and $Y_{u,t}=X_u(t,t+H]$.  Split ownership is defined by
target dates, so validation and test lookbacks may cross the start of their
target interval but their targets may not.

\section{Difficulty and curriculum}
The default user difficulty is the median training-window persistence nMSE:
\[
 d_u=\operatorname{median}_t
 \frac{\operatorname{MSE}(Y_{u,t},X_{u,t,L}\mathbf 1_H)}
      {(\operatorname{sd}(X_{u,t})+\varepsilon)^2}.
\]
Only accessible training targets enter this score. Constant or non-finite users
are removed before ranking by default, and ties use stable source user IDs.

\begin{center}
\begin{tabularx}{\linewidth}{@{}lY@{}}
\toprule
Method & Contract \\
\midrule
\code{uniform} & All users remain equally likely. \\
\code{easy_to_hard} & Cumulatively reveal users in increasing $d_u$. \\
\code{hard_to_easy} & Reverse-order control. \\
\code{random_order} & Seeded random cumulative order. \\
\code{exposure_matched} & All users are available immediately, with fixed
probabilities matching easy-to-hard expected exposure. \\
\bottomrule
\end{tabularx}
\end{center}

For phase $k$, $S_k$ samples and active set $A_k$, the target exposure is
\[
 E_u=\sum_k S_k\frac{\mathbf 1[u\in A_k]}{|A_k|}.
\]
The comparison between easy-to-hard and exposure-matched separates temporal
ordering from marginal reweighting.  The default is five equal-duration
phases, an initial 20\% of users, and linear pacing; square-root and quadratic
pacing are supported.

\section{Models, metrics, and uncertainty}
The standalone forecasting core provides DLinear and PatchTST, global and
instance normalization, MSE and normalized losses, exact optimizer-step
budgets, and query-date sampling.  Primary metrics are global MSE, equal-user
MSE, worst-10\%-user MSE, convergence curves, difficulty-quantile metrics, and
realized versus expected exposure. The reversible instance-normalization type
is named \code{RevIN}, without a redundant ``Normalization'' suffix.

One run seed is the only randomness identifier for that run. It must fix user
ordering, user/window sampling, model initialization, and optimization. Across
seed runs, reported dispersion intentionally includes all of those sources.

\section{Implementation path and provenance}
The scientific treatment is isolated in \code{src/curriculum}: difficulty
scores, schedules, controllers, and curriculum-aware fitting. General tensor
handling follows the TimeTensors interfaces through separate
\code{data/core.py}, \code{sampling.py}, \code{frames.py}, \code{io.py},
\code{splits.py}, \code{statistics.py}, and \code{loaders.py}. Ordinary
construction, fitting, manifests, summaries, and plots remain respectively in
\code{model_loading}, \code{training}, \code{pipeline}, \code{results}, and
\code{visualization}; no sibling checkout is imported.

\noindent The source-adapted external packages are pinned as follows:
\begin{flushleft}\small
PatchTST package: \texttt{external\_models/patchtst}\\
Upstream: \texttt{yuqinie98/PatchTST}\\
Revision: \texttt{204c21efe0b39603ad6e2ca640ef5896646ab1a9}\\[2pt]
DLinear package: \texttt{external\_models/dlinear}\\
Upstream: \texttt{cure-lab/LTSF-Linear}\\
Revision: \texttt{0c113668a3b88c4c4ee586b8c5ec3e539c4de5a6}
\end{flushleft}
The package snapshots are
byte-identical to TimeTensors and expose the common
\code{lags, dim, horizon} tensor boundary.

The executed scientific path is:
\begin{verbatim}
load panel and chronological target splits
compute training-only user difficulty and a stable ordering
construct one fixed-step sampler for the selected curriculum
fit one unchanged model/optimizer across every curriculum phase
evaluate common query dates; aggregate global, user, W10, and quantile metrics
\end{verbatim}

\section{Experiment grid}
\begin{center}
\small
\begin{tabularx}{\linewidth}{@{}lYYYY@{}}
\toprule
Mode & Datasets & Settings & Models & Seeds/budget \\
\midrule
\code{test} & Electricity & $504{:}168$ & PatchTST & seed 1, 20 steps \\
\code{full} & ETTh1, Electricity, Traffic, Solar, Weather, Exchange &
$168{:}24$, $336{:}48$, $504{:}168$ & PatchTST & seeds 1--3 \\
\code{ultra} & full grid & full grid & PatchTST, DLinear & seeds 1--3 \\
\bottomrule
\end{tabularx}
\end{center}
All five curricula are crossed with each selected configuration. ETTh1 uses
all seven non-date variables rather than an OT-only derivative.

\section{Execution and recovery}
The DGX front is \code{curriculum.slurm}; Selena uses
\code{curriculum_selena.slurm}. Both source the same implementation and default
to \code{STAGES=train,tables}. Submit
\begin{verbatim}
EXPERIMENT_MODE=test sbatch curriculum.slurm
EXPERIMENT_MODE=full sbatch curriculum.slurm
EXPERIMENT_MODE=ultra sbatch curriculum.slurm
\end{verbatim}
Replacing the filename by \code{curriculum_selena.slurm} preserves the same
experiment on Selena. Workflow roots \code{LOGS_ROOT} and \code{OUTPUTS_ROOT}
default to \path{logs/} and \path{outputs/}; the Selena front overrides them to
\path{logs_selena/} and \path{outputs_selena/}.
The front sources \code{src/slurm/run_curriculum_experiment.sh}. Training and
table implementations live under \code{src/slurm/}. Resubmission allocates or
resumes current-manifest runs and executes only incomplete seeds. Use
\code{STAGES=train} or \code{STAGES=tables} only for recovery.

Slurm workflows never submit a publisher or execute Git commands. After any
terminal state, including failure, cancellation, or timeout, the user runs
\code{bash publish_job.sh <job-id>} from the project root. The script verifies
\code{main}, sources \code{$HOME/codes/proxy.sh} (or
\code{PROXY_SCRIPT_PATH}), and fast-forward pulls \code{origin/main} before
artifact selection or staging. Lightweight publication selects the exact
local or Selena stdout/stderr pair when a job ID is supplied, all log trees
otherwise, and only aggregate \path{outputs*/reports/} artifacts. Detailed
publication adds all non-binary outputs and diagnostics. Both tiers exclude
\code{*.pt}, \code{*.npy}, and \code{*.cbm}. The script then pushes
\code{origin/main} without opening a pull request. A
non-fast-forward pull stops without a merge commit, and unrelated staged paths
remain outside the publisher commit.

\section{Data and artifact contract}
The loader discovers \code{config.json} beside the selected CSV or at an
explicit path. Shared fields are top-level, project overrides live under
\code{curriculum_learning}, and explicit run settings override both. For
\code{drop_users}, null inherits, an empty list retains every CSV user, and a
nonempty run list replaces the configured default. The project reads CSV
directly and does not create or consume TimeTensors dataset caches.

The ordered relative identity below the selected output root is
\begin{verbatim}
curriculum/dataset/L_H/backbone/method/difficulty/aggregation/
  pacing/phases/initial_fraction/run_n/seed_n/
\end{verbatim}
Every named model config owns one directory. Steps, batch size, learning rate,
evaluation cadence, splits, strides, and quantile count are pipeline configs;
device and Slurm placement are runtime configs. One seed fixes every stochastic
choice in a repetition. Mode selects a subset of these paths and is neither a
path nor signature field.

Each seed saves resolved and dataset configurations, difficulty and curriculum
plans, model state, history, loss plot, global/per-user losses,
\code{results.json}, and \code{all_losses.pt}. The enclosing
\code{manifest.json} uses \code{schema_version: 1}, records configs, optional
plain provenance, launch attempts, per-seed state, required artifacts, and one
of the four overall states \code{not_run}, \code{running}, \code{interrupted},
or \code{completed}. A seed state may additionally be \code{ready}.
Run identity contains only the schema, identity/model configs,
pipeline/experiment configs, and seeds. It never fingerprints or hashes code, Slurm
files, data, weights, logs, outputs, or directories; code/data changes are
manual rerun decisions. The schema changes only for a deliberate global
contract break, and \code{new} creates another repeat with unchanged parameters.

The default \code{overwrite_exact} policy skips an exact completion, resumes an
exact interruption, and allocates the next \code{run_n} for a changed pipeline.
The overall run remains \code{running} with \code{ready_at_utc} while finished
seed states are \code{ready}; completion occurs immediately after that
configuration's producer process returns successfully. Later configuration or
table failures preserve completed runs and interrupt only unfinished work. The completed manifest is authoritative,
without later file hashing or synchronization checks. Tables accept
distinct/latest/average pipeline policies and selected/latest/distinct/average
repeat policies. Explicit pipeline filters select configurations and must match
even for one run; \code{SELECTED_RUNS.txt} stores only automatic or pinned exact
repeats. Reports live under \path{outputs/reports/curriculum/mode/} and record
requested filters and obtained input manifests. Other schema versions are rerun or archived, never read
through a compatibility fallback.

\section{Lightweight validation}
Run the schedule, dataset-config, workflow, Torch curriculum, and smoke tests
under \code{src/tests/}. Do not run publication training locally.
\end{document}
